% =============================================================================
%  Spurious Precision: Why Shift-Share Designs Overstate Monetary Policy
%                      Effects on Local Employment
%  Methods — draft of August 11, 2026
%  Author: Abdallah Dalis
%
%  Written after Results and Discussion, per step 5 of
%  00_Resources/research-paper-writing-process.md.
%
%  Standard for this section: enough detail to reproduce. Every construction
%  below is traced to the line of code that performs it. Where a quantity is
%  not recorded anywhere I could source, it is marked TODO rather than guessed.
% =============================================================================

\section{Data and Methods}

\subsection{Data}

% SOURCE: monetary_policy_labor.R lines 55-86; causal-did scripts section 1.0.
Four sources are used, all public.

County employment comes from the Quarterly Census of Employment and Wages,
assembled as a county-by-quarter panel covering 2002 through 2019. The outcome
is the log of county employment. Observations with non-positive employment are
dropped.

% ⚠️ ADDED August 12, 2026. This is not housekeeping. The panel stores ln_emp as
% 0 where employment is 0, so a screen on the log variable passes every row and
% silently retains 32 county-quarters that double the outcome's standard
% deviation. Stating the screen is on levels is what stops a replicator, or a
% future version of this project, from reproducing that error. See the header of
% Dalis_Abdallah_LocalProjections_BothWindows.R.
The screen is applied to the employment level rather than to its logarithm. The
panel records log employment as zero in the small number of county-quarters
where employment itself is zero, so a finiteness check on the log variable would
retain them. Those observations sit roughly nine log points below the median
county and are concentrated in very small counties and at county
boundary-recodes; they are recording artifacts rather than economic zeros.

% [EXPOSURE-DEPENDENT] ⚠️ August 12, 2026. As currently built this paragraph is
% FALSE: the measure is not the share described, because the numerator sums every
% NAICS hierarchy level while the denominator sums one. The 2,889 and 2,887
% counts below are also downstream of a `drop if share > 1` screen that a correct
% build never triggers. See causal-did/EXPOSURE_DEFECT_2026-08-12.md.
% The paragraph is left as written because it states the INTENDED measure, which
% the rebuild will deliver. Re-verify the counts before submission.
Exposure comes from County Business Patterns. It is the share of county
employment in construction and manufacturing in 2002, the first year of the
panel. The share is fixed at its 2002 value for every quarter, so it cannot
respond to later policy.

% RESOLVED August 11, 2026, by Dalis_Abdallah_ExposureDescriptives.R. All three
% counts floating around the project record are correct; they are different
% objects, which is why none of them could be reconciled by argument:
%   3,228 = raw QCEW panel
%   2,889 = CBP 2002 exposure file
%   2,887 = joined estimation sample (2 counties in CBP have no QCEW match)
% Row attrition: 207,672 joined county-quarters, 196,148 estimated. The 11,524
% difference is the four leads times 2,887 counties, which is 11,548. Close
% enough to be the lead construction and nothing else.
% UPDATED August 12, 2026 for the rebuilt exposure measure. The old figures
% (CBP 2,889 / joined 2,887 / 207,672 / 161,528) were downstream of the defective
% build and of its `drop if share > 1` screen, which deleted 282 counties.
% ✅ FINAL, same day. All counts below are read off the re-run model output, not
% derived by arithmetic here:
%   /tmp -> causal-did/output/tables/run_zlb-full_pos-off.log and run_zlb-excl_pos-off.log
%   joined            224,952 full / 174,968 excl
%   main specification 224,952 full / 174,968 excl   (no leads, so no attrition)
%   lead specification 212,468 full / 162,484 excl
% The lead attrition is 12,484 rows in the full window, which is four leads times
% 3,121 of the 3,127 counties. It is the lead construction and nothing else.
The Quarterly Census of Employment and Wages panel covers 3{,}228 counties. The
County Business Patterns exposure file covers 3{,}127. Joining them, and
requiring an observed funds-rate change, leaves \textbf{3{,}127 counties} over
72 quarters, and that is the estimation sample in both windows.

The 2002 County Business Patterns county file records the District of Columbia
as a statewide entry rather than a county, and it is the only jurisdiction it
treats that way. It is recoded here to its standard county FIPS and retained;
the alternative, dropping every statewide entry, would silently exclude a real
jurisdiction along with the fifty genuine residual records. The District has the
lowest exposure in the country at $0.020$.

Attrition beyond that is mechanical. The joined panel has 224{,}952
county-quarter observations in the full window and 174{,}968 excluding 2009
through 2012. The main specification uses all of them. The specification that
adds four pre-period leads uses 212{,}468 and 162{,}484 respectively, the
difference being rows whose leads are not observed. Counts reported in the paper
are estimation samples unless stated otherwise.

Macroeconomic series come from FRED: the effective federal funds rate, total
nonfarm payrolls, and the consumer price index. The funds rate is averaged
within quarter and first-differenced.

% SOURCE: monetary_policy_labor.R lines 63-71. Sheet name matters.
The identified policy shock is the orthogonalized Bauer--Swanson monetary policy
surprise, taken from the Federal Reserve Bank of San Francisco's published
workbook, sheet ``Monthly (update 2023)''. Monthly surprises are summed within
quarter. The series is constructed to be orthogonal to macroeconomic and
financial information available before each meeting.

\subsection{The shift-share difference-in-differences design}

% SOURCE: causal-did/Dalis_Abdallah_CausalDiD_Analysis.R line 256.
The main specification is
%
\begin{equation}
\ln E_{ct} \;=\; \beta \,\bigl(\Delta f_t \times s_c\bigr) \;+\; \alpha_c \;+\; \tau_t \;+\; \varepsilon_{ct},
\label{eq:did}
\end{equation}
%
where $\ln E_{ct}$ is log employment in county $c$ in quarter $t$, $\Delta f_t$
is the quarterly change in the federal funds rate, and $s_c$ is the 2002
exposure share. County fixed effects $\alpha_c$ absorb permanent differences in
level, and quarter fixed effects $\tau_t$ absorb every national shock, including
$\Delta f_t$ itself.

% NOTE: this is why dffr does not appear as its own regressor. Worth stating
% because a reader will look for it.
The funds rate change therefore does not enter separately. It is collinear with
the quarter fixed effects and is absorbed by them. Only the interaction is
identified, and $\beta$ is a difference in differences: the extra employment
response of a more exposed county relative to a less exposed one.

Standard errors are clustered by county throughout.

Four further specifications are reported. A distributed-lag version adds four
quarterly lags of the interaction. A growth version replaces the outcome with
the quarterly change in log employment. A pre-trends version adds four leads,
and is described below. An event-study version replaces continuous exposure
with an indicator for above-median exposure.

% ⚠️ DECIDED August 11, 2026. Do not silently switch these.
% The event-study parameterization is NOT used for the pre-trends claim. See
% FIGURE_PLAN_2026-08-11.md, "Figure 3: parameterization settled".
The pre-trends test uses the continuous parameterization, not the event-study
one. The two disagree about which leads reject, and the continuous version is
the one on which this paper's randomization inference was computed. Reporting
one test and inferring from the other would be an error, and the choice is
stated here so it is not made by accident later.

\subsection{Local projections}

% SOURCE: monetary_policy_labor.R lines 176-203.
The second design estimates, at each horizon $h=0,\dots,12$,
%
\begin{equation}
100 \times \bigl(\ln E_{c,t+h} - \ln E_{c,t-1}\bigr)
\;=\; \theta_h \,\bigl(z_t \times \tilde{s}_c\bigr) \;+\; \text{(county and quarter means)} \;+\; u_{ct},
\label{eq:lp}
\end{equation}
%
where $z_t$ is the policy shock and $\tilde{s}_c$ is exposure standardized to
mean zero and unit variance. Both the outcome and the interaction are passed
through a two-way within transformation, removing county and quarter means,
before estimation without an intercept.

The specification is estimated twice at every horizon. The first pass uses the
realized funds rate change as $z_t$. The second substitutes the Bauer--Swanson
surprise. Nothing else changes: same panel, same transformation, same horizons.
That is the point of the exercise.

% ⚠️ SCALE WARNING. A reader will otherwise try to compare 0.022 against 0.198.
Two features of this design make its coefficients incomparable to
equation~\eqref{eq:did}, and both are easy to miss. The outcome is scaled by
100 and cumulated from $t-1$ to $t+h$, and exposure is standardized rather than
left in share units. So $\theta_h$ is a percentage response per standard
deviation of exposure, while $\beta$ is a log-point response per unit of
exposure share. \textbf{The two coefficients are not on the same scale and
should never be compared directly.} The comparison this paper makes is between
two local projections that share a scale, not between the projections and the
difference in differences.

% SOURCE: monetary_policy_labor.R line 197.
Standard errors in the local projections are clustered by quarter rather than
by county. This is the appropriate choice: the shock has no cross-sectional
variation, so the residual dependence that matters runs across counties within
a quarter. It also means inference rests on 72 clusters rather than on the
number of observations. Section~\ref{sec:limits} returns to this.

\subsection{Randomization inference}

% SOURCE: causal-did/Dalis_Abdallah_RandomizationInference.R, seed 20260811.
The pre-trends result is checked without relying on the clustered variance
estimator. Exposure is permuted across counties, the model is re-estimated on
the permuted assignment, and the exercise is repeated 200 times per
construction. The spread of the placebo estimates is the null distribution, and
the reported $p$-value is the share of permutations whose estimate exceeds the
true one in absolute value. The seed is fixed at 20260811.

% ⚠️ Both schemes are reported, and the second must NOT be read as stronger
% evidence. See PIPELINE_RECONCILIATION.md section 4f.
Two permutation schemes are run. The first shuffles exposure freely across all
counties. The second shuffles only within state. The within-state scheme was
intended to preserve the geographic concentration of industrial structure that
free permutation destroys. It does not do what was intended: restricting the
shuffle to same-state counties means exposure can only move to a county with
similar exposure, which makes the randomization less aggressive and narrows the
null mechanically. Its $p$-values are therefore anti-conservative and are not
quoted as evidence. Both are reported so the reader can see why.

\subsection{Sample windows}

% ⚠️ Honest statement of a decision made without a remembered rationale.
Every result is reported for two samples: the full 2002 to 2019 panel, and the
same panel excluding 2009 through 2012.

The earlier version of this work excluded those years. The zero lower bound is a
defensible reason to do so, since the funds rate carries little information
about the policy stance when it is pinned near zero. It is not the recorded
reason, because no rationale was recorded at the time.

Rather than reconstruct one, both windows are reported throughout. A
justification written after seeing the results is worth nothing, and showing
both leaves no choice to defend. Where the two disagree, the disagreement is
reported. Where they do not, that is stated as well.

\subsection{Software and reproducibility}

% SOURCE: MEMORY.md, "Stata license lapsed, August 2026. R is the only
% pipeline." Do NOT describe the .do files as runnable.
All estimates are produced in R. Fixed-effects models use \texttt{fixest}, with
\texttt{haven} for Stata-format inputs, \texttt{sandwich} for clustered variance
estimation in the local projections, and \texttt{readxl} for the Bauer--Swanson
workbook.

An earlier version of the difference-in-differences pipeline was written in
Stata. It is retained for provenance only and cannot be executed. Where the two
pipelines disagreed, the disagreement was traced rather than reconciled by
choosing the preferred answer, and the R implementation is canonical. The
resulting reconciliation is documented in full, including the finding that a
routine choice about how lagged variables are constructed cut standard errors by
40 to 45 percent without meaningfully moving the point estimates.

% This sentence is the honest version of a limitation, not a boast. Keep it.
That last observation deserves emphasis, because it is this paper's own thesis
appearing inside its own code: a defensible-looking construction choice, of a
kind no published table would ever display, manufactured precision that was not
there.
